% !TeX TS-program = xelatex
% ============================================================================
% 郭睢阳 - 个人简历
% 基于 adongwanai/LLM-Resume-Template (CC BY 4.0)
% 原始仓库: https://github.com/adongwanai/LLM-Resume-Template
% ============================================================================
\documentclass{resume-photo}

% 字体（FandolSong 随 TeX Live 分发，无需额外安装；如已装 Noto 可改回 Noto Serif CJK SC）
\usepackage{xeCJK}
\setCJKmainfont{FandolSong-Regular.otf}[BoldFont=FandolSong-Bold.otf, ItalicFont=FandolKai-Regular.otf]
\usepackage{fontspec}

% 重点高亮（关键技术/量化成果标红加粗，提升可扫描性）
\usepackage{xcolor}
\definecolor{hlred}{RGB}{192,0,0}
\newcommand{\hl}[1]{\textbf{\textcolor{hlred}{#1}}}

% 紧凑排版
\usepackage{enumitem}
\setlist[itemize]{itemsep=2pt, topsep=2pt}
\newcommand{\ExpItem}[3]{%
  \noindent\textbf{#1｜#2} \hfill {\footnotesize #3}\\[-2pt]
}

% ====== 个人信息 ======
\ResumeName{郭睢阳}
\ResumePhoto{photo.jpg}  % 替换为你的照片文件名，不想要头像就注释掉

\begin{document}

\ResumeContacts{
  15893346463,%
  \ResumeUrl{mailto:1327217178@qq.com}{1327217178@qq.com},%
  \textnormal{西北工业大学 | 控制工程 · 硕士 | 2024.09—2027.06}%
}

\ResumeTitle

% ====================== 简介 ======================
\noindent 西北工业大学控制工程专业硕士在读（2027年6月毕业），方向机器人与人工智能。具备从机械设计、电路设计到嵌入式软件和上位机开发的\hl{全流程研发能力}。现于 IDEA 研究院：视启未来从事 \hl{VLA（视觉-语言-动作模型）与世界模型}算法研发，推动双臂机器人柔性操作策略的\hl{训练、实时推理优化与真机部署}。熟悉 ROS、Linux 应用开发与深度学习，参与\hl{国家级重点课题}，累计 \hl{6 项专利（4 发明 / 2 实用新型）}与 \hl{2 项软件著作权}。

% ====================== 教育 ======================
\section{教育背景}
\ResumeItem{西北工业大学\hl{（985）}}
[\textnormal{控制工程 | 硕士 | 机器人与人工智能方向 | GPA: 89.56/100（前 30\%）}]
[2024.09—2027.06]

\ResumeItem{郑州大学\hl{（211）}}
[\textnormal{自动化（卓越班）| 学士 | 机器人与人工智能方向 | GPA: 3.54/4.0（前 20\%）}]
[2020.09—2024.06]

% ====================== 技能 ======================
\section{专业技能}
\begin{itemize}
  \item \textbf{AI / 具身智能}: \hl{VLA}（\hl{\(\pi\)0.5}、\hl{X-VLA}）训练与\hl{真机部署}、\hl{flow-matching policy}、action chunking、\hl{RTC 实时推理}、\hl{里程碑世界模型}；\hl{PyTorch（DDP）}/ \hl{JAX-Flax（FSDP）}分布式训练
  \item \textbf{深度学习}: PyTorch / JAX 训练全流程、\hl{DINOv3 / SigLIP} 视觉表征、CNN、语音识别；\hl{Triton kernel}、\hl{CUDA Graphs} 推理加速
  \item \textbf{编程语言}: \hl{C/C++}（精通）、\hl{Python}（精通）、MATLAB（熟练），熟悉 Linux 驱动/应用开发
  \item \textbf{机器人}: \hl{ROS/ROS2}（精通），\hl{导航、SLAM、运动规划、逆运动学求解}、点云处理
  \item \textbf{嵌入式 / 硬件}: \hl{STM32/51/Arduino}（精通，IMX6ULL / STM32F4）、\hl{FreeRTOS}；PCB 设计（Altium Designer、嘉立创）、3D 打印
  \item \textbf{工具}: \hl{Vibecoding}（Claude Code 等 AI 辅助编程）、Git、PyQt5 上位机开发、Inventor 3D 建模
\end{itemize}

% ====================== 项目经历 ======================
\section{项目经历}

\ResumeItem{串并联混合轮式仿人机器人}
[\textnormal{\hl{国家级重点课题} | 关键省重点实验室项目成员}]
[2024.09—至今]

\begin{itemize}
  \item \textbf{项目简介}: 设计搭建串并联混合型轮式仿人机器人系统，包含 AGV 小车、并联机构、两个 6 自由度机械臂（自研 6 轴/JAKA）及 2 自由度头部，实现农业采摘、实验教学、工业场景应用
  \item \textbf{技术栈}: Ubuntu 20.04、\hl{ROS Noetic}、\hl{STM32F4}、点云处理、路径规划、\hl{神通数据库}
  \item \textbf{工作内容}: 编写机器人总控程序，实现\hl{三维地图构建、路径规划、目标识别和动态空间避障}；实现点云发布、碰撞检测、\hl{逆运动学求解}；设计自研多自由度机械臂控制方案，实现\hl{六自由度机械臂控制频率提升 100\%}；基于 STM32F4 编写下位机控制程序；搭建神通数据库实现高并发数据增删改查
\end{itemize}

\ResumeItem{AI 大模型智能语音交互机器人}
[\textnormal{\hl{校优秀毕业设计} | 独立开发者}]
[2023.12—2024.09]

\begin{itemize}
  \item \textbf{项目简介}: 结合前沿 NLP 技术、\hl{AI 大模型}和 AI 对口型技术，实现智能家用语音聊天机器人
  \item \textbf{技术栈}: 树莓派 4B、Python、\hl{PyQt5}、\hl{PyTorch}、ROS、FreeRTOS
  \item \textbf{工作内容}: 设计 AI 大模型智能交互软件，实现\hl{录音、语音识别、大模型交互}；实现语音合成、\hl{AI 数字人对口型}功能；实现智能家居控制交互；获得 \hl{2 项软件著作权}
\end{itemize}

\ResumeItem{中医理疗仪系列产品}
[\textnormal{\hl{挑战杯省级二等奖} | 核心研发成员}]
[2021.03—2023.09]

\begin{itemize}
  \item \textbf{三代产品}: 智能葫芦灸理疗仪（一代）→ 新型智能化温针理疗仪（二代）→ 家用可移动灸罐理疗仪（三代）
  \item \textbf{技术栈}: \hl{STM32}、\hl{FreeRTOS}、\hl{模糊 PID}、传感器融合、蓝牙通信
  \item \textbf{工作内容}: 主导三代理疗仪\hl{全流程研发}（机械设计→电路设计→嵌入式软件→上位机），完成 STM32 嵌入式系统开发，负责 STM32 芯片外围电路\hl{PCB 设计}，温度/图像/压力等多种传感器数据接入，伺服与步进电机驱动开发
\end{itemize}

% ====================== 实习经历 ======================
\section{实习经历}

\ResumeItem{IDEA 研究院：视启未来}
[\textnormal{VLA 与世界模型算法实习生 · 具身智能团队}]
[2026.04—至今]

\begin{itemize}
  \item 以开源 \hl{\(\pi\)0.5}（openpi）为基座训练双臂机器人叠衣策略，经数据精选与两阶段初始化做出\hl{新 SOTA 配方}，动作误差 MAE@50 较原 SOTA \hl{降低约 77\%}，并完成\hl{真机部署上机}
  \item 主导 \hl{VLA 推理实时化}：实现 \hl{RTC 实时分块推理}与 \hl{Triton kernel 优化}，单步推理耗时由约 256ms 优化至 32ms（\hl{约 8 倍}），末端抖动\hl{降低 81\%}
  \item 实现 \hl{AWBC / RECAP advantage} 离线策略改进流水线，开展 \hl{cross-embodiment 跨本体混训}与权重空间模型合并等系统实验
  \item 完成清华 \hl{X-VLA} 架构在自有双臂平台真机 bring-up，用固件级 IK 将推理周期由 \hl{1.4s 压至 0.33s}；建立离线视觉消融门禁，\hl{定位并修复模型“视觉失明”根因缺陷}
  \item 自研\hl{里程碑世界模型（LMVLA）}，零训练挖掘 milestone 子目标并注入 \(\pi\)0.5 动作专家，在 \hl{LIBERO / RoboTwin} 仿真基准上系统评测
\end{itemize}

\ResumeItem{中秦禾瑞（陕西）科技有限公司}
[\textnormal{导航与建图算法工程师（兼职）}]
[2025.11—2026.01]

\begin{itemize}
  \item 负责果园智能巡检小车\hl{三维建图和导航算法}开发，设备选型、调试与传感器数据处理
  \item 接入\hl{雷达原始数据实现 3D 建图}，基于三维地图处理为\hl{栅格地图}进行巡检小车导航
\end{itemize}

\ResumeItem{华为数字能源实践营}
[\textnormal{学员（培训）}]
[2025.07]

\begin{itemize}
  \item 学习数字能源云微服务 \hl{DDD 架构}与 AI 工程化部署，掌握\hl{提示词工程}与 DevOps 部署流程
\end{itemize}

\ResumeItem{东风汽车有限公司}
[\textnormal{参观学习}]
[2023.08]

\begin{itemize}
  \item 参观学习汽车汽油/柴油发动机内部结构和工作原理，了解汽车全过程工程制造流程
\end{itemize}

% ====================== 竞赛与荣誉 ======================
\section{竞赛与荣誉}
\begin{itemize}
  \item 第十六届挑战杯河南省大学生课外学术科技作品竞赛 \hfill \hl{省级二等奖（2023）}
  \item 第二届高校电气电子工程师创新大赛 \hfill \hl{华中赛区一等奖（2023）}
  \item 第十六届西门子杯智能制造挑战赛 \hfill 西部赛区二等奖（2022）
  \item 蓝桥杯嵌入式设计与开发组 \hfill 省级三等奖（2022）
  \item 哈尔滨工业大学郑州研究院特训营 \hfill 优秀营员（2022）
  \item 郑州大学二等奖学金、三好学生 \hfill 2022—2023, 2024—2025
\end{itemize}

% ====================== 专利 ======================
\section{专利与证书}
\begin{itemize}
  \item \hl{发明专利 4 项}、\hl{实用新型专利 2 项}、\hl{软件著作权 2 项}
  \item 英语六级（CET-6）: 473 分
\end{itemize}

\end{document}
